% ============================================================
%  Conclusion.tex  –  Game Tracker
% ============================================================

\chapter{Conclusión}

\section{Resultados obtenidos}

El desarrollo de Game Tracker permitió construir un sistema funcional de
extremo a extremo: desde el diseño de la base de datos relacional hasta la
interfaz móvil que el usuario opera en su dispositivo. El catálogo de
videojuegos puede gestionarse completamente a través de la aplicación,
incluyendo el registro de juegos con su estado de progreso, la escritura de
notas con soporte de respuestas anidadas y la posibilidad de reaccionar a
las publicaciones de otros usuarios.

Desde el punto de vista técnico, el proyecto cumplió con el objetivo de
aplicar patrones de diseño de forma deliberada y justificada, no como un
requisito superficial, sino como decisiones de arquitectura que resolvieron
problemas concretos durante el desarrollo.

\section{Reflexión sobre los Patrones de Diseño}

Trabajar con patrones de diseño durante este proyecto cambió la forma en que
se tomaron las decisiones de implementación. Antes de aplicarlos, la tendencia
natural era escribir toda la lógica donde pareciera más conveniente en el
momento, lo que terminaba generando controladores muy extensos y difíciles de
seguir. Los patrones obligaron a hacerse preguntas concretas: \textit{¿quién
es responsable de este comportamiento? ¿esta lógica puede cambiar en el futuro?
¿quién necesita saber sobre esto?}

El patrón que más impacto tuvo en la calidad del código fue \textbf{Strategy},
especialmente en la gestión de las estrategias de eliminación de notas. Lo que
comenzó como dos casos manejados con \texttt{if-else} creció hasta tres
variantes con reglas distintas según el rol del usuario, y la estructura del
patrón absorbió ese crecimiento sin necesidad de reescribir la lógica
existente.f

El patrón \textbf{Singleton} en \texttt{AuthService} resolvió un problema real
de coordinación: sin él, mantener el estado de sesión coherente entre múltiples
pantallas habría requerido pasar el token como argumento a través de cada widget,
ensuciando el código con información que no le corresponde a la mayoría de
las clases.

\textbf{ChangeNotifier} demostró que la separación entre vista y lógica no es
solo una buena práctica teórica: al aislar toda la lógica de reacciones en
\texttt{ReactionTimelineController}, fue posible modificar el comportamiento del
modal sin tocar la pantalla principal, y actualizar la interfaz sin volver a
hacer llamadas HTTP innecesarias.

Finalmente, \textbf{MVC} como base estructural del backend facilitó la
incorporación progresiva de nuevas funcionalidades ---reacciones, moderación,
eliminación en cascada--- sin que los cambios en una capa afectaran a las demás.

\section{Lecciones aprendidas}

\begin{itemize}
    \item La arquitectura cliente-servidor con una API REST como contrato entre 
    ambos lados permitió desarrollar el backend y el frontend de forma paralela 
    e independiente, lo cual simplificó el proceso de trabajo.

    \item El despliegue en la nube expuso problemas reales que no aparecen en 
    entornos locales, como la configuración correcta de CORS, la gestión de 
    variables de entorno y los tiempos de respuesta de una base de datos remota.

    \item Definir bien el sistema de roles desde el inicio del proyecto evitó 
    tener que refactorizar la lógica de permisos más adelante, cuando ya había 
    código que dependía de ella.

    \item Centralizar las reglas de autorización en un único servicio 
    (\texttt{NoteAuthorizationService}) facilitó la escritura de pruebas 
    unitarias, ya que toda la lógica de permisos vive en un solo lugar.
\end{itemize}

\section{Trabajo a futuro}

Como extensiones naturales del sistema se identifican las siguientes mejoras:

\begin{itemize}
    \item Integración con APIs externas ---como RAWG o IGDB--- para importar 
    automáticamente la información de los juegos, incluyendo portada, género 
    y descripción oficial.
    \item Notificaciones en tiempo real cuando otro usuario reacciona a una 
    nota propia, utilizando WebSockets o eventos enviados por el servidor.
    \item Un panel de administración web dedicado para gestionar el catálogo 
    sin necesidad de usar la aplicación móvil.
    \item Ampliación del sistema de roles para incluir listas de juegos 
    privadas y la posibilidad de compartirlas con usuarios específicos.
\end{itemize}